\begin{theorem}[deriv-conj-conj]
\label{deriv_conj_conj'}
\lean{deriv_conj_conj'}
% at /home/runner/work/trough/trough/library/PrimeNumberTheoremAnd/ZetaConj.lean:12.0-28.6
\leanok
Let $f : \mathbb{C} \to \mathbb{C}$ be a function at $p \in \mathbb{C}$ with derivative $a$.
    Then the derivative of the function $g(z) = \overline{f(\overline{z})}$ at $\overline{p}$ is
    $\overline{a}$.
\end{theorem}
\begin{proof}
\leanok
We proceed by case analysis on whether $f$ is differentiable at $p$. If $f$ is differentiable
    at $p$, then we can apply the previous theorem. If $f$ is not differentiable at $p$, then
    neither is $g$, and both derivatives have the default value of zero.
\end{proof}
